\section{Results}
\label{sec:results}

\subsection{A 1.6M-parameter model against foundation models}
\label{sec:main}

Table~\ref{tab:main} compares CroFreMo trained from scratch against the
strongest reproduced baseline on each corpus; the full matrix of fifteen
baselines is in the appendix. Every entry is a mean over three seeds unless
marked. The model wins on the three corpora whose label is a paroxysmal
event or pattern with adequate positives (TUSZ, CHB-MIT, IIIC) and on TUEP,
ties the strongest foundation model on TUEV, TUAB, and ADFD, and loses on
the two sleep-staging corpora, on CAUEEG, on TUAR, and on Siena. The
largest margins are on seizure detection, where the strongest baselines are
not the largest models: on CHB-MIT, CBraMod trained from scratch reaches
$0.317 \pm 0.167$ and its released pretrained weights $0.233 \pm 0.080$,
against $0.699 \pm 0.045$ for a model twenty times smaller.

\begin{table}[t]
\centering
\small
\caption{CroFreMo (1.6M parameters, trained from scratch) against the
strongest reproduced baseline per corpus. Three seeds, mean $\pm$ std; $(1)$
marks a single seed. Pretrained rows are marked ``pre''. The pretrained
CroFreMo row uses the fixed recipe of Section~\ref{sec:setup}.}
\label{tab:main}
\begin{tabular}{llllll}
\toprule
Corpus & Metric & Strongest baseline & & CroFreMo & CroFreMo pre \\
\midrule
TUSZ & AUC-PR & FFCL & $0.545 \pm 0.024$ & $\mathbf{0.639 \pm 0.029}$ & $0.653 \pm 0.033$ \\
CHB-MIT & AUC-PR & TFM-Tokenizer pre & $0.627 \pm 0.021$ & $\mathbf{0.699 \pm 0.045}$ & $0.699\,(1)$ \\
IIIC & $\kappa$ & REVE pre & $0.436 \pm 0.002$ & $\mathbf{0.479 \pm 0.008}$ & $0.434\,(1)$ \\
TUEP & AUROC & EEGPT scratch & $0.786 \pm 0.018$ & $\mathbf{0.810 \pm 0.003}$ & $0.769\,(1)$ \\
TUEV & $\kappa$ & REVE pre & $0.685 \pm 0.032$ & $0.690 \pm 0.029$ & $0.642\,(1)$ \\
TUAB & bal.\ acc. & ST-Transformer & $0.820 \pm 0.004$ & $0.816\,(1)$ & $0.805\,(1)$ \\
ADFD & bal.\ acc. & BIOT scratch & $0.525 \pm 0.017$ & $0.505 \pm 0.050$ & $0.453\,(1)$ \\
CAUEEG & bal.\ acc. & BIOT scratch & $0.561 \pm 0.009$ & $0.525 \pm 0.012$ & $0.503\,(1)$ \\
Siena & AUC-PR & REVE pre & $0.518 \pm 0.096$ & $0.170 \pm 0.070$ & $0.456 \pm 0.034$ \\
Sleep-EDF & $\kappa$ & ContraWR & $0.692 \pm 0.012$ & $0.642 \pm 0.011$ & $0.675\,(1)$ \\
ISRUC & $\kappa$ & CBraMod pre & $0.754 \pm 0.006$ & $0.702 \pm 0.007$ & $0.672\,(1)$ \\
TUAR & $\kappa$ & CBraMod pre & $0.715\,(1)$ & $0.620 \pm 0.030$ & $0.592\,(1)$ \\
\bottomrule
\end{tabular}
\end{table}

\paragraph{Where the margin comes from.}
The seizure-detection margin is not the coupling token: the same model with
the coupling pathway removed (raw waveform tokens only) reaches
$0.671 \pm 0.026$ on TUSZ and $0.667 \pm 0.047$ on CHB-MIT
(Table~\ref{tab:coupling}). Two architectural controls locate it instead
(Table~\ref{tab:arch}). Collapsing the learned filterbank to a single
broadband band costs $0.11$ on TUSZ and $0.42$ on CHB-MIT; replacing the
tri-axial encoder by a flat encoder that mixes electrodes before attention
costs $0.22$ and $0.53$. The band-decomposed token grid and the factorized
attention over it are what let a small model learn rare-positive seizure
windows, and the baselines that sit at the level of these ablations are
precisely the ones whose tokens or encoders lack one of the two.

\begin{table}[t]
\centering
\small
\caption{Architectural controls, raw-token frontend, three seeds. Removing
either the band decomposition or the tri-axial factorization collapses
seizure detection.}
\label{tab:arch}
\begin{tabular}{lccc}
\toprule
 & 8 bands, tri-axial & 1 band, tri-axial & 8 bands, flat (0.9M) \\
\midrule
TUSZ (AUC-PR) & $0.671 \pm 0.026$ & $0.558 \pm 0.041$ & $0.448 \pm 0.048$ \\
CHB-MIT (AUC-PR) & $0.667 \pm 0.047$ & $0.245 \pm 0.068$ & $0.136 \pm 0.016$ \\
\bottomrule
\end{tabular}
\end{table}

\subsection{Coupling as token content}
\label{sec:coupling-results}

Table~\ref{tab:coupling} isolates the coupling token: the same encoder, the
same recipe, with the tokenizer producing either waveform tokens only or the
duplex grid of Section~\ref{sec:tokenizer}. On TUEV the coupling token is
decisive, $+0.15$ in $\kappa$ over three seeds, four standard deviations of
the raw-token model. On IIIC the gain is small but consistent
($+0.013$, three seeds). On the ten remaining corpora it is zero within
seed variance, including the two seizure-detection corpora on which the
model wins by the largest margins. Coupling content is therefore not what
makes the model competitive in general; it is what makes it competitive on
epileptiform event classification.

\begin{table}[t]
\centering
\small
\caption{Coupling as token content: raw waveform tokens vs.\ the duplex
grid, same encoder and recipe, three seeds each.}
\label{tab:coupling}
\begin{tabular}{llccc}
\toprule
Corpus & Metric & Raw tokens & Duplex & $\Delta$ \\
\midrule
TUEV & $\kappa$ & $0.536 \pm 0.034$ & $0.690 \pm 0.029$ & $\mathbf{+0.15}$ \\
IIIC & $\kappa$ & $0.466 \pm 0.008$ & $0.479 \pm 0.008$ & $+0.01$ \\
CHB-MIT & AUC-PR & $0.667 \pm 0.047$ & $0.699 \pm 0.045$ & $+0.03$ \\
TUSZ & AUC-PR & $0.671 \pm 0.026$ & $0.639 \pm 0.029$ & $-0.03$ \\
TUEP & AUROC & $0.810 \pm 0.001$ & $0.810 \pm 0.003$ & $0$ \\
ADFD & bal.\ acc. & $0.501 \pm 0.040$ & $0.505 \pm 0.050$ & $0$ \\
CAUEEG & bal.\ acc. & $0.517 \pm 0.027$ & $0.525 \pm 0.012$ & $0$ \\
TUAR & $\kappa$ & $0.623 \pm 0.007$ & $0.620 \pm 0.030$ & $0$ \\
Sleep-EDF & $\kappa$ & $0.650 \pm 0.002$ & $0.642 \pm 0.011$ & $0$ \\
ISRUC & $\kappa$ & $0.701 \pm 0.025$ & $0.702 \pm 0.007$ & $0$ \\
Siena & AUC-PR & $0.171 \pm 0.130$ & $0.170 \pm 0.070$ & $0$ \\
\bottomrule
\end{tabular}
\end{table}

\paragraph{The gain sits on periodic epileptiform discharges.}
Table~\ref{tab:perclass} breaks the TUEV gain down by class. It concentrates
on generalized and lateralized periodic discharges (GPED $+0.27$, PLED
$+0.09$ in F1), the two event types defined by a slow rhythm on which fast
activity rides, which is the signature the modulation index of
Eq.~(\ref{eq:mvl}) measures. The coupling model gives up eye movements
($-0.25$), a slow deflection with no cross-frequency structure, and neither
model learns the rare spike-and-sharp-wave class. This is the pattern the
mechanism predicts: the token helps exactly where the label is a
cross-frequency morphology.

\begin{table}[t]
\centering
\small
\caption{TUEV per-class F1, raw tokens vs.\ duplex, averaged over three
seeds each (29,421 test windows).}
\label{tab:perclass}
\begin{tabular}{lcccccc}
\toprule
 & SPSW & GPED & PLED & EYEM & ARTF & BCKG \\
\midrule
Raw tokens & 0.16 & 0.55 & 0.52 & 0.57 & 0.40 & 0.88 \\
Duplex & 0.00 & $\mathbf{0.82}$ & $\mathbf{0.61}$ & 0.31 & 0.56 & 0.92 \\
\bottomrule
\end{tabular}
\end{table}

\subsection{The token inside another encoder}
\label{sec:transplant}

To separate the tokenizer from the encoder it was designed with, we
replace CBraMod's patch embedding by CroFreMo's tokenizer, keeping CBraMod's
criss-cross encoder, positional encoding, classifier and published recipe
unchanged; every (electrode, token row) becomes one CBraMod channel and rows
are pooled after the encoder, so the head is identical
(Table~\ref{tab:transplant}). On TUEV the transplant beats CBraMod's own
tokenizer by $0.068$, and switching the coupling off inside the same
transplant gives back $0.037$ of it. On TUSZ the transplant ties the native
tokenizer and the coupling on/off gap is $0.052$. In both encoders,
therefore, the coupling content is what the gain is made of; whether the
gain shows up against the native tokenizer depends on how much of the
cross-frequency structure the encoder can already recover from band
tokens. The tri-axial encoder, which attends along a frequency axis,
recovers it on seizure windows and not on periodic discharges; CBraMod,
which has no such axis, benefits on both. A mean-pooled transplant that
collapses the band axis before CBraMod's encoder, reported in the appendix,
loses on most corpora, which is the same statement in the negative.

\begin{table}[t]
\centering
\small
\caption{CroFreMo's tokenizer inside CBraMod's encoder. CBraMod's own
tokenizer is a three-seed mean; the transplants are single seeds.}
\label{tab:transplant}
\begin{tabular}{lccc}
\toprule
 & CBraMod tokenizer & Transplant, coupling off & Transplant, coupling on \\
\midrule
TUEV ($\kappa$) & $0.564 \pm 0.019$ & $0.595$ & $\mathbf{0.632}$ \\
TUSZ (AUC-PR) & $0.482 \pm 0.043$ & $0.422$ & $0.475$ \\
\bottomrule
\end{tabular}
\end{table}

\subsection{Pretraining, reported to convergence}
\label{sec:pretrain-results}

The last column of Table~\ref{tab:main} is the pretrained model finetuned
under one fixed recipe. It improves on the from-scratch model where labels
are scarce or noisy (Siena, where positives are rare; TUSZ; Sleep-EDF) and
loses where labels are plentiful (TUEV, ADFD, TUAR), a pattern the analysis
in the appendix explains rather than excuses: lr-matched frozen probes show
that the pretrained representation is informative (on TUEV, $0.65$ against
$0.39$ for a random encoder), label-fraction curves show the harm growing
with the number of labels, and clean-transfer controls at the pretraining
patch length reproduce the TUEV loss with the entire backbone loaded. The
coupling reconstruction term of Eq.~(\ref{eq:objective}) helps the frozen
representation on TUEV and hurts it on TUSZ; band-normalized amplitude
reconstruction alone \citep{yu2026understanding} is neither better nor
worse in finetuning. We report these as findings about a small model with a
strong inductive bias in the clinical label regime, not as a property of
the objective.
